% ============================================================
% OP_GUT2_repclass_recon_v3.tex  (PROPOSAL-ONLY; v3 per GPT review
% of v2). Changes vs v2: SU(3)^3 SM-sector line rewritten
% (mandatory microfix); N_{<=4} chirality clarification; S_2
% restricted to nonzero charges with the Y=0 colour-pair caveat;
% D_0 renamed to minimal type-count and fenced; the "<=6
% multiplets" classification bound replaced by a bounded
% first-pass audit framing; NEW: coloured weak-singlet sector
% computed exactly (Lemma C_odd and Lemma C_4) -- the four-
% multiplet case admits a genuine chiral anomaly-free family,
% exhibited with an explicit witness, which upgrades the
% no-overclaim statement from a caution to a demonstrated fact;
% insertion scoping per review (V full, N_{<=4}/S_2 short form,
% D_0 proposal-only, no long proof sketches in the article).
% NO preprint edits.
% ============================================================
\section*{OP-GUT2 conditional solution-space audit, v3
(proposal-only)}

\textit{Status target: candidate Level~A algebra inside the
restricted class (for the sectors actually computed); Level~B
conditional as ECT application; OP-GUT2 itself remains Open.}

\subsection*{1. Conventions}
All fermion fields are left-handed Weyl fields; $u^c,d^c,e^c$
(and, in its branch, $\nu^c$) denote left-handed conjugates of the
right-handed Standard-Model fields.
$A(\mathbf 3)=+1$, $A(\bar{\mathbf 3})=-1$;
$\mathbf 2$ of $SU(2)$ is pseudo-real, so
$(\mathbf 1,\mathbf 2)_Y\oplus(\mathbf 1,\mathbf 2)_{-Y}$ is
vector-like.
A vector-like pair is $(R_c,R_w,Y)\oplus(\bar R_c,R_w,-Y)$.

\subsection*{2. Definition: P7-restricted admissible matter class}
Left-handed Weyl fields valued only in tensor products of the
already assumed structures: the colour module $E$ or $E^*$
($\mathbf 3$ or $\bar{\mathbf 3}$), the weak doublet/singlet
structure, and a compact-$U(1)$ charge-lattice label.
This is a structural restriction defining a test class, not a
derivation of matter representations.
The exclusion of $\mathbf 6,\mathbf 8,\dots$ and higher weak
isospin is not derived from anomaly cancellation and not derived
from ECT dynamics; it is a minimality/available-module assumption
defining the test class.

\subsection*{3. Minimality (classification convention)}
At most one copy of each chiral type $(R_c,R_w,Y)$; no vector-like
pairs counted as part of the minimal chiral core; no sterile
singlets unless an explicit $\nu^c$ branch is being analysed.
This is a classification convention, not a dynamical selection
rule.

\subsection*{4. Anomaly system on the class}
As in v2: $[SU(3)]^3$, $[SU(3)]^2U(1)$, $[SU(2)]^2U(1)$,
$\mathrm{grav}^2U(1)$, $[U(1)]^3$, and the colour-weighted Witten
parity $\sum_{\rm doublets} d_c\equiv0\ (\mathrm{mod}\ 2)$.

\subsection*{5. SM-pattern sector (computed; reuses recorded
AN3/$N_c$ algebra)}
Pattern $\mathcal P_{\rm SM}=
\{(\mathbf 3,\mathbf 2)_{Y_Q},(\bar{\mathbf 3},\mathbf 1)_{Y_u},
(\bar{\mathbf 3},\mathbf 1)_{Y_d},(\mathbf 1,\mathbf 2)_{Y_L},
(\mathbf 1,\mathbf 1)_{Y_e}\}$.
$[SU(3)]^3$: $Q$ contributes two colour triplets ($d_w=2$),
$u^c$ and $d^c$ one anti-triplet each, so
\[
2A(\mathbf 3)+A(\bar{\mathbf 3})+A(\bar{\mathbf 3})=2-1-1=0 .
\]
Witten: $3+1=4$, even.
Remaining conditions and their full solution structure: as
recorded (AN3 cubic, two hypercharge branches, degenerate
$Y_Q=0$ branch, normalisation freedom; $N_c=3$ case of the
recorded $N_c$-classification).
$\nu^c$ branches: (a)~none (baseline); (b)~$Y_{\nu^c}=0$, all
conditions unchanged; (c)~free $Y_{\nu^c}$, opening the recorded
$B-L$ flat direction.

\subsection*{6. Structural lemmas (own algebra)}
\paragraph{Lemma V (vector-like transparency).}
A vector-like pair contributes zero to all five local conditions
and an even number to the Witten count; anomaly conditions do not
constrain vector-like additions.
\paragraph{Lemma N$_{\le4}$ (small pure-singlet sector).}
There is no anomaly-free chiral set consisting solely of $n\le4$
fields $(\mathbf 1,\mathbf 1)_{y_i}$ with all $y_i\neq0$.
Here ``chiral'' means after removing vector-like $(y,-y)$ pairs
and neutral sterile singlets $y=0$, consistently with the
minimal-core convention.
Proof as in v2 ($n=3$: $\sum y_i=0\Rightarrow\sum y_i^3=3y_1y_2y_3$;
$n=4$: symmetric-function argument forces $s(y_1y_2-y_3y_4)=0$,
both branches vector-like).
\paragraph{Lemma S$_2$ (small coloured weak-singlet sector).}
With no weak doublets and at most two nonzero coloured
weak-singlet multiplets, $[SU(3)]^3$ forces one $\mathbf 3$ and
one $\bar{\mathbf 3}$, and $[SU(3)]^2U(1)$ forces
$Y_{\bar{\mathbf 3}}=-Y_{\mathbf 3}$: a vector-like pair.
If the coloured charges vanish, the pair
$(\mathbf 3,\mathbf 1)_0\oplus(\bar{\mathbf 3},\mathbf 1)_0$ is
still vector-like under colour.
\paragraph{Lemma D$_0$ (colourless doublet sector at minimal
type-count).}
As in v2 (pseudo-reality plus $[SU(2)]^2U(1)$ plus Witten parity
close the sector at minimal type-count).
This is a minimal-type-count statement, not a classification of
arbitrary multi-copy colourless doublet sectors.
[Per review: D$_0$ stays proposal-only or one short sentence in
any eventual insertion.]

\subsection*{7. Coloured weak-singlet sector: exact audit
(new in v3)}
\paragraph{Lemma C$_{\rm odd}$.}
With no weak doublets, $[SU(3)]^3$ requires equal total numbers of
$\mathbf 3$ and $\bar{\mathbf 3}$ multiplets ($d_w=1$ throughout);
hence no coloured weak-singlet sector with an odd number of
coloured multiplets closes in this restricted class.
In particular the three-multiplet case is excluded.
\paragraph{Lemma C$_4$ (four-multiplet case, solved).}
Take $(\mathbf 3,\mathbf 1)_{a_1},(\mathbf 3,\mathbf 1)_{a_2},
(\bar{\mathbf 3},\mathbf 1)_{b_1},(\bar{\mathbf 3},\mathbf 1)_{b_2}$,
all charges nonzero.
The conditions reduce to $a_1{+}a_2{+}b_1{+}b_2=0$ (from
$[SU(3)]^2U(1)$; $\mathrm{grav}$ is then automatic) and
$a_1^3{+}a_2^3{+}b_1^3{+}b_2^3=0$.
With $s:=a_1{+}a_2=-(b_1{+}b_2)$ the cubic condition gives
$3s\,(b_1b_2-a_1a_2)=0$.
\emph{Branch $s\neq0$:} $b_1b_2=a_1a_2$ and $b_1{+}b_2=-(a_1{+}a_2)$
force $\{b_1,b_2\}=\{-a_1,-a_2\}$: two vector-like pairs.
\emph{Branch $s=0$:} $a_2=-a_1$, $b_2=-b_1$, and all five
conditions hold identically; for $b_1\neq\pm a_1$ no vector-like
pair exists, and no gauge-invariant mass bilinear exists
($\mathbf 3_y\otimes\bar{\mathbf 3}_{y'}$ requires $y{+}y'=0$,
excluded; $\mathbf 3\otimes\mathbf 3$ and
$\bar{\mathbf 3}\otimes\bar{\mathbf 3}$ contain no singlet; the
$\epsilon$-contraction needs three triplets).
\emph{Conclusion:} the four-multiplet coloured weak-singlet sector
admits a genuine two-parameter chiral anomaly-free family
\[
\{(\mathbf 3,\mathbf 1)_{a},(\mathbf 3,\mathbf 1)_{-a},
(\bar{\mathbf 3},\mathbf 1)_{b},(\bar{\mathbf 3},\mathbf 1)_{-b}\},
\qquad a,b\neq0,\ b\neq\pm a,
\]
e.g.\ $a=1,b=2$: colour-anomaly-free, $U(1)$-protected chiral
quartets, sterile under $SU(2)$.
Direct sums of anomaly-free sets are anomaly-free, so such
quartets could in principle accompany the SM pattern; they are
excluded from the minimal core only by the minimality convention.
\paragraph{Scope note.}
Mixed sectors (coloured quartets plus extra neutral singlets,
larger doublet-bearing patterns) remain open and are not guessed.

\subsection*{8. Bounded first-pass audit (reframed scope)}
The article-facing claim is a \emph{bounded first-pass audit}: the
SM-pattern core, the optional $\nu^c$ branch, and the small
non-SM sectors controlled by the lemmas above (V, N$_{\le4}$,
S$_2$, D$_0$, C$_{\rm odd}$, C$_4$); larger chiral patterns remain
open.
This is an audit, not a classification of all P7-restricted matter
contents; no fixed multiplet-count bound is presented as a
classification scope.

\subsection*{9. No-overclaim lemma (updated; now with witness)}
Anomaly cancellation, Witten parity, and lattice quantisation do
not by themselves select the Standard-Model matter content; this
is not merely a caution but a computed fact: the restricted class
admits non-SM chiral anomaly-free content (the C$_4$ quartet
family) alongside the SM-pattern solution branches.
Within the class, the SM pattern is distinguished by the
minimality convention and by carrying the full
colour{+}weak{+}hypercharge chiral structure, not by anomaly
consistency alone; identifying it as the realised matter content
requires additional input, ultimately the open dynamical
derivation OP-GUT2.

\subsection*{10. Insertion scoping (per review; for the eventual
block, not executed)}
Lemma~V in full; N$_{\le4}$ and S$_2$ in short form; D$_0$
proposal-only or one sentence; C$_{\rm odd}$/C$_4$ proposed in
short form with the explicit family displayed once; no long proof
sketches in the article.
Summary sentence of the small-sector algebra, per review draft:
``Small-sector algebra closes vector-like additions, pure singlets
up to four nonzero charges, and at most two coloured
weak-singlets; the four-multiplet coloured weak-singlet sector
admits an explicit chiral family; larger sectors remain open.''

\subsection*{11. Questions for GPT}
(Q1)~Approve Lemma C$_{\rm odd}$ and Lemma C$_4$ (algebra above)?
(Q2)~C$_4$ family in the eventual article block: short form with
the displayed family (my preference, since it gives the
no-overclaim lemma a witness), or proposal-only?
(Q3)~Audit framing of \S8 acceptable as the article-facing scope?
(Q4)~If approved: exact insertion blocks next (D1-style), placed
in the colour-completion/OP-GUT2 discussion, with status
candidate-Level-A-inside-class / Level-B-ECT / OP-GUT2-Open?
